\documentclass[11pt]{amsart}
\usepackage[margin=1.1in]{geometry}
\usepackage{amssymb,amsmath,amsthm}
\usepackage{booktabs}
\usepackage{hyperref}

\newtheorem{theorem}{Theorem}[section]
\newtheorem{lemma}[theorem]{Lemma}
\newtheorem{proposition}[theorem]{Proposition}
\newtheorem{corollary}[theorem]{Corollary}
\newtheorem*{mainq}{Main Question}
\theoremstyle{definition}
\newtheorem{definition}[theorem]{Definition}
\newtheorem{hypothesis}[theorem]{Hypothesis}
\newtheorem{remark}[theorem]{Remark}

\newcommand{\F}{\mathbb{F}_3}
\newcommand{\Z}{\mathbb{Z}}
\newcommand{\C}{\mathbb{C}}
\newcommand{\R}{\mathbb{R}}
\newcommand{\cN}{\mathcal N}
\newcommand{\Index}{\operatorname{Index}}
\newcommand{\Aut}{\operatorname{Aut}}
\newcommand{\Mon}{\operatorname{Mon}}
\newcommand{\Fer}{\operatorname{Fer}}
\newcommand{\ch}{\operatorname{ch}}
\newcommand{\str}{\operatorname{str}}
\newcommand{\Tr}{\operatorname{Tr}}
\newcommand{\res}{\operatorname{res}}
\newcommand{\maj}{\operatorname{maj}}
\newcommand{\gw}{\operatorname{gw}}
\newcommand{\Reg}{\mathrm{Reg}}
\newcommand{\SH}{\widetilde{\mathrm{SH}}{}^3}
\newcommand{\rSH}{\mathrm{rSH}^3}
\newcommand{\Hc}{\mathrm{H}}
\newcommand{\orb}{/\!\!/}

\title{Holomorphic $\mathcal N=1$ vertex operator superalgebras of minimal index}
\author{[Authors to be added]}
\date{Draft, 10 October 2026}

\begin{document}

\begin{abstract}
Gaiotto and Johnson-Freyd asked whether for every $k$ there is a holomorphic $\cN=1$ vertex operator
superalgebra of central charge $12k$ whose supersymmetric index is $24/\gcd(k,24)$, the smallest value
compatible with topological modular forms. They answered the question for $k\le5$ by orbifolding the
lattice superalgebra of a self-dual ternary code of index $24$ by a cyclic group of code automorphisms.
Using the codes of a companion paper we answer it for every $k$, within their orbifold framework.
When $8\nmid k$ their cyclic construction applies. When $8\mid k$ no cyclic group can work, and we
orbifold instead by the quaternion group $Q_8$ ($k\equiv8,16\pmod{24}$) or by $Q_8\times\Z/3$
($k\equiv0\pmod{24}$). The main new input is the 't~Hooft anomaly of a quaternion group of code
automorphisms: it vanishes if and only if $8\mid k$. Its restrictions to cyclic subgroups vanish already
for $4\mid k$; one class of restricted supercohomology is invisible to all of them, and we detect it
through the $(-1)$-twisted sector, by showing that the character of a would-be orbifold fails to be
$S$-invariant when $k\equiv4\pmod8$. We also show that the orbifold has index $\Index(C)/|G|$ for every
group $G$ of code automorphisms whose abelian subgroups are cyclic, whose only involution is $-1$, and
whose elements of odd order form a central subgroup of order at most $3$ that fixes no coordinate and as
many even as odd maximal codewords. For $k=24$ this gives a holomorphic
$\cN=1$ vertex operator superalgebra of central charge $288$ and index $1$.
\end{abstract}

\maketitle

\section{Introduction}\label{sec:intro}

A self-dual ternary code $C\subseteq\F^n$ determines the odd unimodular lattice
$\Lambda(C)=\{a/\sqrt3: a\in\Z^n,\ a\bmod3\in C\}$ and hence a holomorphic vertex operator
superalgebra $V_{\Lambda(C)}$ of central charge $n$, which Gaiotto and Johnson-Freyd equip with an
$\cN=1$ structure \cite[\S2]{GJF}. Its supersymmetric (Ramond--Ramond) index is the index of the code
\cite[Prop.~2.4]{GJF},
\[
  \Index(C)=\#\{\text{even maximal codewords}\}-\#\{\text{odd maximal codewords}\},
\]
where a codeword is maximal if none of its entries is $0$, and even or odd according to the parity of
the number of its entries equal to $1$. Under the Stolz--Teichner conjecture a holomorphic $\cN=1$
vertex operator superalgebra of central charge $12k$ defines a class in $\pi_{24k}\mathrm{TMF}$ whose
image in $\mathrm{MF}_{12k}$ is its index times $\Delta^k$ \cite[\S5]{GJF}, and the smallest multiple
of $\Delta^k$ in the image of $\pi_{24k}\mathrm{TMF}\to\mathrm{MF}_{12k}$ is
$\frac{24}{\gcd(k,24)}\Delta^k$ \cite[Prop.~4.6]{Hopkins}. Gaiotto and Johnson-Freyd ask whether these
values are attained:

\begin{mainq}[{\cite[\S0, \S5]{GJF}}]
For each $k$, does there exist a holomorphic $\cN=1$ VOA with central charge $c=12k$ and
supersymmetric index $\frac{24}{\gcd(k,24)}$?
\end{mainq}

The index is defined up to sign, the parity of the Ramond sector being a convention
\cite[Rem.~5.3]{GJF}; ``index $m$'' means index $\pm m$ throughout. Gaiotto and Johnson-Freyd answered
the question for $k\le5$ with the following theorem, found their codes by computer search, and asked for
a systematic construction \cite[Question~0.1, Thm.~5.6, Ex.~5.7]{GJF}.

\begin{quote}
\textbf{Theorem} (\cite[Thm.~5.6]{GJF}). Let $C$ be self-dual of length $12k$ with $\Index(C)=24$,
put $m=\gcd(24,k)$, and let $g\in\Aut(C)$ have order $m$. Suppose that (1) if $m$ is even,
$g^{m/2}=-1$, and (2) if $3\mid m$, $g^{m/3}$ fixes as many even as odd maximal codewords and permutes
the coordinates in $3$-cycles only. Then $V_{\Lambda(C)}\orb\langle g\rangle$ has index $24/m$.
\end{quote}

In a companion paper \cite{I} we constructed explicit self-dual codes of index $24$ in every length
$12k$, with the automorphisms that this theorem and the present paper need, and we showed that for
$8\mid k$ no cyclic group can work: an automorphism $g$ with $g^4=-1$ forces $48\mid\Index(C)$
\cite[Thm.~D]{I}. Here we supply the vertex-algebra side.

\begin{theorem}[Main Theorem]\label{thm:main}
Assume Hypothesis~ORB (\S\ref{sec:orb}) and the cited inputs (A1)--(A10) of \S\ref{sec:setting}. For
every $k\ge1$ there is a holomorphic $\cN=1$ vertex operator superalgebra of central charge $12k$ and
index $24/\gcd(k,24)$, namely $V_{\Lambda(C)}\orb G$ for a code $C$ and a group $G\subseteq\Aut(C)$ as in
Table~\ref{tab:main}.
\end{theorem}

\begin{table}[ht]
\centering
\begin{tabular}{@{}cllcll@{}}
\toprule
$\gcd(k,24)$ & $k\bmod 24$ & index & $G$ & code & seed of the code \\
\midrule
$1$ & $\pm1,\pm5,\pm7,\pm11$ & $24$ & $1$ & \cite[Thm.~B]{I} & Golay \\
$2$ & $\pm2,\pm10$ & $12$ & $\langle-1\rangle$ & \cite[Thm.~B]{I} & Golay \\
$3$ & $\pm3,\pm9$ & $8$ & $\langle h\rangle\cong\Z/3$ & \cite[Thm.~C(a)]{I} & $E_{36}$ of \cite{GJF} \\
$4$ & $\pm4$ & $6$ & $\langle J\rangle\cong\Z/4$ & \cite[Thm.~C(b)]{I} & Golay \\
$6$ & $\pm6$ & $4$ & $\langle -h\rangle\cong\Z/6$ & \cite[Thm.~C(a)]{I} & $E_{36}$ of \cite{GJF} \\
$8$ & $\pm8$ & $3$ & $Q_8=\langle J,K\rangle$ & \cite[Thm.~C(b)]{I} & Golay \\
$12$ & $12$ & $2$ & $\langle Jh\rangle\cong\Z/12$ & \cite[Thm.~C(c)]{I} & $E_{72}$ of \cite[App.~A.2]{I} \\
$24$ & $0$ & $1$ & $Q_8\times\langle h\rangle$ & \cite[Thm.~C(c)]{I} & $E_{72}$ of \cite[App.~A.2]{I} \\
\bottomrule
\end{tabular}
\medskip
\caption{The orbifolds of Theorem~\ref{thm:main}. Here $J^2=K^2=-1$, $KJ=-JK$, and $h$ has order $3$, no
fixed coordinate, a vanishing twined count $T_h(C)=0$, and commutes with $J$ and $K$.}
\label{tab:main}
\end{table}

For $8\nmid k$ this is \cite[Thm.~5.6]{GJF} applied to the codes of \cite{I}, as noted there.
Sections~\ref{sec:cyclic} and~\ref{sec:index} reprove both halves of that theorem, the vanishing of the
anomaly and the value of the index, from the same inputs and in a form that also covers the non-cyclic
groups. For $8\mid k$ the new ingredients are the following three theorems. All automorphisms of codes
act on $V_{\Lambda(C)}$ through the lift of \cite[Thm.~3.1]{GJF}.

\begin{theorem}\label{thm:Q8}
Assume (A1)--(A10) and Hypothesis~ORB. Let $C$ be a self-dual ternary code of length $12k$, and let $J,K\in\Aut(C)$ satisfy $J^2=K^2=-1$ and
$KJ=-JK$. The action of $Q_8=\langle J,K\rangle$ on $V_{\Lambda(C)}$ is non-anomalous if and only if
$8\mid k$.
\end{theorem}

\begin{theorem}\label{thm:Q8Z3}
If moreover $h\in\Aut(C)$ has order $3$, no fixed coordinate, and commutes with $J$ and $K$, then
$Q_8\times\langle h\rangle$ acts non-anomalously on $V_{\Lambda(C)}$ if and only if $24\mid k$.
\end{theorem}

\begin{theorem}\label{thm:index}
Let $C$ be self-dual of length $12k$ and $G\subseteq\Aut(C)$ a subgroup such that
\begin{enumerate}
\item[(i)] every abelian subgroup of $G$ is cyclic;
\item[(ii)] if $|G|$ is even, then $-1\in G$ is the only involution of $G$;
\item[(iii)] the elements of odd order form a central subgroup of order $1$ or $3$, and in the second case
  it is generated by an $h$ without fixed coordinate such that $T_h(C)=0$.
\end{enumerate}
If $G$ acts non-anomalously on $V_{\Lambda(C)}$, then, under Hypothesis~ORB, every orbifold
$W=V_{\Lambda(C)}\orb G$ is a holomorphic $\cN=1$ vertex operator superalgebra with
$\Index(W)=\pm\Index(C)/|G|$.
\end{theorem}

Here $T_h(C)=\sum_{x\in C,\,hx=x}f(x)$ is the signed count of $h$-fixed maximal codewords
(\S\ref{sec:codes}); it vanishes iff $h$ fixes as many even as odd maximal codewords. All groups of
Table~\ref{tab:main} satisfy (i)--(iii).

\subsection*{The idea for \texorpdfstring{$Q_8$}{Q8}}
The anomaly of a $Q_8$-action lies in the restricted supercohomology group $\rSH(Q_8)$, of order $32$.
The Frame-shape computation of \cite[Ex.~4.4]{GJF} gives its restrictions to the cyclic subgroups, which
vanish iff $4\mid k$. A count (Proposition~\ref{prop:kernel}) shows that exactly one nonzero class
$\kappa$ restricts to zero on every cyclic subgroup. Since every abelian subgroup of $Q_8$ is cyclic, no
twisted--twined partition function of a commuting pair can see $\kappa$. We detect it through the
$(-1)$-twisted sector. As a projective $Q_8$-module this sector does not depend on $C$: it is the
spinor module of the Clifford algebra of the frame's zero modes (Lemma~\ref{lem:minus}). So the
character that a $Q_8$-orbifold would have can be written down, and three eta identities show that it
is not $S$-invariant when $k\equiv4\pmod8$ (Proposition~\ref{prop:Sdefect}). For $8\mid k$, additivity
of anomalies under direct sums of codes finishes the proof. For the record, the anomaly is $3k\cdot a$,
where $a$ is the anomaly of eight Majorana fermions in the real regular representation of $Q_8$, a class
of order $8$ that takes the value $\tfrac58$ on $S^3/Q_8$ (Remark~\ref{rem:freefermions}).

\subsection*{What the results rest on}
The fermionic case of the orbifold theorem \cite[Thm.~4.2]{GJF} is only sketched in \cite{GJF}; for
bosonic vertex operator algebras it is a theorem \cite{Kirillov,DNR,McRae}. Hypothesis~ORB
(\S\ref{sec:orb}) states what we use from it: existence of the orbifold for a non-anomalous group, and
the decomposition of its NS and Ramond sectors into invariants of twisted sectors, with genuine actions
of the centralisers. We also use the description of the Majorana layer and the multiplier principle for
cyclic groups from \cite[\S4]{GJF}, which are stated there without proof, and standard theorems on lattice
vertex superalgebras and their twisted modules \cite{BK,Li,DZ,DRY}. Gaiotto and Johnson-Freyd's own
answers for $k\le5$ rest on the same inputs. On the code side we use \cite{I}; for $12\mid k$ one
numerical input enters there, the value $24$ of a block-weighted sum over a code of length $72$, which
was computed by three independent programs.

\subsection*{Two remarks}
Holomorphic $\cN=1$ vertex operator superalgebras are closed under tensor products, and for central
charges divisible by $12$ the index is multiplicative. So the theory of central charge $288$ and index
$1$ in Theorem~\ref{thm:main} reduces the Main Question to $k\le24$ (Remark~\ref{rem:reduction}). Under
the Stolz--Teichner conjecture that theory represents a unit of $\pi_{576}\mathrm{TMF}$
(Remark~\ref{rem:TMF}).

\subsection*{Organisation}
\S\ref{sec:setting} collects the setting and the cited inputs. \S\ref{sec:cyclic} treats cyclic groups,
\S\ref{sec:Q8} proves Theorem~\ref{thm:Q8}, \S\ref{sec:Q8Z3} Theorem~\ref{thm:Q8Z3}, and
\S\ref{sec:index} Theorem~\ref{thm:index}. \S\ref{sec:proof} proves the Main Theorem and adds remarks.
The appendix lists the computations that check the explicit formulas.

\section{Setting}\label{sec:setting}

\subsection{Codes}\label{sec:codes}
We follow \cite[\S2]{I}. Let $\varphi\colon\F\to\{0,\pm1\}$ be $\varphi(0)=0$, $\varphi(1)=-1$,
$\varphi(2)=1$, let $f(x)=\prod_i\varphi(x_i)$, and $F(X)=\sum_{x\in X}f(x)$ for $X\subseteq\F^n$. Then
$\Index(C)=F(C)$. The group $\Mon(n)=2^n{:}S_n$ acts on $\F^n$ and on $\R^n$ by $ge_i=s_ie_{p(i)}$ for
$g=(p,s)$, and $\Aut(C)=\{g:gC=C\}$ contains $-1$. Put $N_-(g)=\#\{i:s_i=-1\}$ and
$\varepsilon(g)=\prod_is_i$; then $\varepsilon$ is a character and $f(gx)=\varepsilon(g)f(x)$.

A cycle of $p$ of length $\ell$ whose signs have product $\sigma$ is a \emph{signed cycle}; it
contributes the factor $1-x^\ell$ ($\sigma=1$) or $1+x^\ell=(1-x^{2\ell})/(1-x^\ell)$ ($\sigma=-1$) to
$\det(1-xg)=\prod_k(1-x^k)^{d_k}$. The formal product $\prod_kk^{d_k}$ is the \emph{Frame shape} of $g$,
and we put $\eta_g(\tau)=\prod_k\eta(k\tau)^{d_k}$ and
\[
  \rho_g=\tfrac14\sum_j\theta_j(1-\theta_j)=\tfrac1{24}\sum_kd_k\bigl(k-\tfrac1k\bigr),
\]
where $e^{2\pi i\theta_j}$, $\theta_j\in[0,1)$, are the eigenvalues of $g$ on $\R^n$. The vector space
$\R^n$ has no nonzero $g$-fixed vector iff every signed cycle has $\sigma=-1$, iff $\sum_kd_k=0$. We say
that $g$ has \emph{no fixed coordinate} if $p$ has no fixed point; this is weaker (an element of order $3$
without fixed coordinate fixes the vectors that are constant on its cycles). On a signed $\ell$-cycle,
$g^j$ is diagonal iff $\ell\mid j$, and then it is $\sigma^{j/\ell}$. Hence
\begin{equation}\label{eq:cycles}
  g^j=-1\iff\text{every signed cycle has length }\ell\mid j\text{ with }j/\ell\text{ odd and }\sigma=-1.
\end{equation}
For example $J^2=-1$ means that $J$ consists of $n/2$ signed $2$-cycles with $\sigma=-1$, so
$N_-(J)=n/2$, $J$ has no fixed vector, and its Frame shape is $2^{-n/2}4^{n/2}$.

\subsection{The vertex superalgebra of a code}
Let $C$ be self-dual of length $n=12k$, $\Lambda=\Lambda(C)$ and $V=V_\Lambda$. The \emph{frame} is the
subalgebra $V_{\sqrt3\Z^n}=V_{\sqrt3\Z}^{\otimes n}\subseteq V$. The vector $\chi=\sqrt3(1,\dots,1)$ is
characteristic: $\langle\chi,\lambda\rangle\equiv|\lambda|^2\pmod2$ for $\lambda\in\Lambda$. Let $\sigma$
be the parity automorphism of $V$.

\begin{proposition}[{\cite[\S2, Prop.~2.4]{GJF}; \cite{BK,Li,DZ}}]\label{prop:VC}
\begin{enumerate}
\item $V$ is a strongly rational holomorphic vertex operator superalgebra of central charge $n$.
\item $\tau=\tfrac1{\sqrt3}\sum_i\bigl(e^{\sqrt3e_i}+e^{-\sqrt3e_i}\bigr)$, for a suitable cocycle, is an
  $\cN=1$ supercurrent: $\tau_{(0)}\tau=2\omega$, $\tau_{(1)}\tau=0$, $\tau_{(2)}\tau=\tfrac{2n}3\mathbf 1$.
\item The Ramond sector $V_R$, the unique $\sigma$-stable irreducible $\sigma$-twisted module, is
  $V_{\Lambda+\chi/2}$, with parity $(-1)^{\langle\chi,\lambda\rangle}$ on $M(1)\otimes e^{\lambda+\chi/2}$.
\item The states of $V_R$ of weight $c/24=n/24$ are the vectors $e^{\nu_w}$, one for each maximal
  codeword $w$, with $\nu_w=\iota(w)/(2\sqrt3)$ and $\iota\colon\F\to\{0,1,-1\}$; the parity of
  $e^{\nu_w}$ is $f(w)$.
\item On every $g$-twisted $V$-module with $g\tau=-\tau$ the modes $G_r=\tau_{r+1/2}$, $r\in\Z$, satisfy
  $G_0^2=L_0-c/24$. In particular $\str_{V_R}q^{L_0-c/24}=\Index(C)$.
\end{enumerate}
\end{proposition}

\begin{proof}[Comments]
(1) is \cite[\S2]{GJF} with \cite[Thm.~4.1]{BK} for odd lattices; (3) is \cite[Prop.~5.4]{Li} applied
to $\chi/2$, with uniqueness from \cite[Thm.~7.7]{DZ}. (4): a Ramond momentum is
$\nu=(2a+3\cdot\mathbf 1)/(2\sqrt3)$ with $a\in\Z^n$, $a\bmod3\in C$, so
$|\nu|^2/2=\sum_i(2a_i+3)^2/24\ge n/24$, with equality iff every $a_i\in\{-1,-2\}$. Then $w=-(a\bmod3)$ is
maximal, $\nu=\nu_w$, and the parity $(-1)^{\sum a_i}$ is $(-1)^{\#\{i:w_i=1\}}=f(w)$. (5) is the
twisted commutator formula \cite[(M4)]{Li} with (2); if $G_0$ is odd, invertible on each weight space
of weight $\ne c/24$, and commutes with an operator $A$, then $A$ has supertrace $0$ there. So
supertraces of operators commuting with $G_0$ only see weight $c/24$, and (4) gives the last claim.
\end{proof}

\subsection{Automorphisms}
By \cite[Thm.~3.1]{GJF} the extension $\widehat\Lambda[2].\Aut(C)$ of $\Aut(C)$ inside $\Aut(V)$
splits. We write $s\colon\Aut(C)\to\Aut(V)$ for the splitting that restricts to the standard action of
$\Mon(n)$ on the frame, $e^{\pm\sqrt3e_i}\mapsto e^{\pm s_i\sqrt3e_{p(i)}}$. It is unique: the restriction
$\operatorname{Hom}(\Lambda,\pm1)\to\operatorname{Hom}(\sqrt3\Z^n,\pm1)$ is bijective, since
$\sqrt3\Z^n$ has odd index in $\Lambda$. Consequently $s_{C\oplus C'}=s_C\otimes s_{C'}$ on
$V_{\Lambda(C\oplus C')}=V_{\Lambda(C)}\otimes V_{\Lambda(C')}$. Every $s(g)$ fixes $\tau$, so $s(G)$
acts by even automorphisms, faithfully, with $s(G)\cap\langle\sigma\rangle=1$; and $s(g)$ maps
$M(1)\otimes e^\lambda$ to $M(1)\otimes e^{g\lambda}$ with a sign \cite[Prop.~4.1]{BK}. We usually write
$g$ for $s(g)$.

\begin{lemma}\label{lem:majnumber}
For $g\in\Aut(C)$, $\langle\chi,(1-g)\chi\rangle=6N_-(g)$, and $s(g)$ maps the summand of $V_R$ with
momentum $\lambda+\chi/2$ to the one with momentum $g(\lambda+\chi/2)$, multiplying the parity by
$(-1)^{N_-(g)}=\varepsilon(g)$.
\end{lemma}

\begin{proof}
$\langle\chi,\chi\rangle=3n$ and $\langle\chi,g\chi\rangle=3\sum_is_i=3(n-2N_-(g))$. The relative parity
is $\tfrac12\langle\chi,(1-g)\chi\rangle=3N_-(g)$ modulo $2$ \cite[Ex.~4.4]{GJF}.
\end{proof}

If $g\in O(\Lambda)$ has no fixed vector, any two lifts of $g$ to $\widehat\Lambda.O(\Lambda)$ are
conjugate by an element of the torus $\widehat\Lambda=\operatorname{Hom}(\Lambda,U(1))$
\cite[Ex.~4.4]{GJF}: if $t\in\widehat\Lambda$, then $t\hat gt^{-1}=\hat g\cdot(t\circ(g-1))$, and
$t\mapsto t\circ(g-1)$ is onto because $g-1$ is injective on $\Lambda$ and $U(1)$ is divisible. In
particular every lift of such a $g$ has the same anomaly as $s(g)$, and $\hat g\sigma$ is conjugate to
$\hat g$, since $\sigma=e^{2\pi i(\chi/2)(0)}$ lies in the torus.

\subsection{Twisted modules}\label{sec:twisted}
For an automorphism $g$ of finite order of a strongly rational holomorphic vertex operator
superalgebra there is a unique $\sigma$-stable irreducible $g$-twisted module $V(g)$
\cite[Thm.~7.7]{DZ}, \cite[Prop.~3.4]{DRY}. The \emph{$g$-twisted Ramond sector} is $V_{R,g}:=V(g\sigma)$.
An \emph{$x$-intertwiner} on a $g$-twisted module $M$, for $x$ commuting with $g$, is an invertible $A$
with $AY_M(v,z)A^{-1}=Y_M(xv,z)$; on $V(g)$ it exists and is unique up to a scalar when $V(g)$ is of
$\C$-type, and then it is even or odd.

\begin{proposition}[{\cite[Thm.~4.2, Rem.~4.4]{BK}, \cite{Li}}]\label{prop:fpf}
Let $g\in\Aut(C)$ have no fixed vector and $\varepsilon(g)=1$. Then $V(g)$ is of $\C$-type,
$V(g)\cong V(g\sigma)$ as graded spaces, and $V(g)=S[g]\otimes T_g$, where $S[g]$ is the $g$-twisted Fock
space of the Heisenberg algebra and $T_g$ is a space of ground states of dimension
$|\det(1-g)|^{1/2}=\prod_kk^{d_k/2}$. The lowest weight is $\rho_g$, and
$\ch V(g)=|\det(1-g)|^{1/2}\prod_k\eta(\tau/k)^{-d_k}$. Moreover
$Z(1,g):=\operatorname{tr}_Vg\,q^{L_0-c/24}=1/\eta_g(\tau)$.
\end{proposition}

\begin{proof}[Comments]
Only $\lambda=0$ contributes to $Z(1,g)$, and the trace on the Fock space is $\prod_m\det(1-gq^m)^{-1}$.
The number of irreducible $g$-twisted modules is $1$ or $2$ according as $\tfrac12\langle\chi,
(1-g)\chi\rangle$ is even or odd \cite[Thm.~4.2]{BK}, i.e.\ according to $\varepsilon(g)$
(Lemma~\ref{lem:majnumber}); the defect formula of \cite{BK} gives $\dim T_g$. The vacuum energy of
$S[g]$ is $\rho_g$ \cite[Thm.~1.7.3]{FLM}, \cite{Lep,DL}; equivalently, $\ch V(g)$ is the
$S$-transform of $Z(1,g)$ \cite[Thm.~7.1]{DZ}, whose leading exponent is $\rho_g-n/24$. Finally
$V(g\sigma)\cong V(g)$ because $g\sigma$ is conjugate to $g$ in the torus.
\end{proof}

\subsection{Anomalies}\label{sec:anomalies}
We use the following statements from \cite{GJF} and the literature it cites.
\begin{itemize}
\item[(A1)] \cite[Thm.~4.2]{GJF}. For a finite group $G$ acting faithfully by even automorphisms on a
  holomorphic vertex operator superalgebra $V$ with $V^G$ regular there is an anomaly
  $\alpha(V,G)\in\SH(G)$, in reduced extended supercohomology. Orbifolds $V\orb G$ exist iff
  $\alpha=0$, and then correspond to trivialisations. For solvable $G$, regularity of $V^G$ follows from
  \cite{CM} applied to the even part of $V$ \cite[\S4]{GJF}.
\item[(A2)] The anomaly is natural: $\alpha(V,H)=\res^G_H\alpha(V,G)$ for $H\subseteq G$, and it is
  transported along isomorphisms of groups.
\item[(A3)] Conjugate actions have equal anomalies.
\item[(A4)] Anomalies add under tensor products: $\alpha(V\otimes V',\rho\otimes\rho')=\alpha(V,\rho)+
  \alpha(V',\rho')$ \cite[\S5]{GJF17}.
\item[(A5)] There are exact sequences, natural in $G$,
  \begin{gather*}
    \Hc^3(G;U(1))\xrightarrow{\iota_1}\rSH(G)\xrightarrow{\gw}\Hc^2(G;\Z/2)\to\Hc^4(G;U(1)),\\
    \rSH(G)\xrightarrow{\iota_2}\SH(G)\xrightarrow{\maj}\Hc^1(G;\Z/2),
  \end{gather*}
  with $\iota_1$ and $\iota_2$ injective \cite[\S4]{GJF}, \cite{WG}, \cite[\S5.4]{GJF17}. The maps
  $\maj$ and $\gw$ are the Majorana and Gu--Wen layers.
\item[(A6)] $\maj(\alpha)(g)=0$ iff an operator on $V_R$ implementing $g$ is even \cite[\S4]{GJF}. For
  the free fermions $\Fer(\rho)$ in an orthogonal representation $\rho$ this gives $\maj=w_1(\rho)=\det
  \rho$.
\item[(A7)] (Multiplier principle \cite[\S4]{GJF}.) Let $G=\langle g\rangle\cong\Z/m$ act with
  vanishing Majorana layer, and put $L=m$ for $m$ even, $L=2m$ for $m$ odd. Then $Z(1,g)/Z(1,1)$ is an
  eigenfunction of $ST^LS^{-1}$, and its eigenvalue $e_m(\alpha)$ depends only on $\alpha\in\rSH(\Z/m)$,
  through a homomorphism $e_m\colon\rSH(\Z/m)\to U(1)$.
\end{itemize}
For $G=\Z/m$ let $\nu_{m,r}\in\rSH(\Z/m)$ be the anomaly of the action on $\Fer(2)=V_{\Z}$ in which the
generator rotates the two fermions through $2\pi r/m$, and $\nu_m=\nu_{m,1}$. Put $m'=2m$ for $m$ even
and $m'=m$ for $m$ odd.

\begin{lemma}\label{lem:calib}
\begin{enumerate}
\item $|\rSH(\Z/m)|=m'$, $|\rSH(Q_8)|=32$, and $\rSH(\Z/m)=\SH(\Z/m)$ for $m$ odd.
\item $e_m(\nu_{m,r})=e^{2\pi ir^2/m'}$. Hence $e_m$ is injective, $\rSH(\Z/m)=\langle\nu_m\rangle\cong
  \Z/m'$, and $\nu_{m,r}=r^2\nu_m$.
\end{enumerate}
\end{lemma}

\begin{proof}
(1) For $G$ cyclic or $Q_8$, $\Hc^3(G;U(1))\cong\Hc^4(G;\Z)\cong\Z/|G|$ and $\Hc^4(G;U(1))\cong
\Hc^5(G;\Z)=0$ (periodic cohomology \cite[XII.7]{CE}, \cite[VI.9]{Brown}), so (A5) gives
$0\to\Z/|G|\to\rSH(G)\to\Hc^2(G;\Z/2)\to0$; and $\Hc^2(\Z/m;\Z/2)$ is $\Z/2$ or $0$, $\Hc^2(Q_8;\Z/2)\cong
(\Z/2)^2$, $\Hc^1(\Z/m;\Z/2)=0$ for $m$ odd. (2) The characters are $Z(1,1)=\sum_pq^{p^2/2}/\eta$ and
$Z(1,g)=\sum_p\zeta^{pr}q^{p^2/2}/\eta$ with $\zeta=e^{2\pi i/m}$. Poisson summation gives
$(Z(1,g)/Z(1,1))|S=\sum_pq^{(p+r/m)^2/2}/\sum_pq^{p^2/2}$, and $T^L$ multiplies $q^{(p+r/m)^2/2}$ by
$e^{2\pi ir^2/m'}$ and fixes the denominator. As $e_m$ is a homomorphism, $\nu_m$ has order at least
$m'=|\rSH(\Z/m)|$.
\end{proof}

\subsection{The orbifold}\label{sec:orb}
For a non-anomalous $G$, \cite[Thm.~4.2]{GJF} constructs the orbifold from a trivialisation, and the
proof of \cite[Thm.~5.6]{GJF} computes its Ramond--Ramond partition function as a sum over twisted
sectors. For bosonic holomorphic vertex operator algebras the corresponding statements are theorems
\cite{Kirillov,DNR,McRae}; in the super case their ingredients exist (\cite[Thms.~4.15,
4.17]{McRae} for the $G\times\langle\sigma\rangle$-crossed structure, \cite{CM} for regularity), but
their assembly is only sketched in \cite{GJF}. We state what we use as a hypothesis.

\begin{hypothesis}[ORB]\label{hyp:orb}
Let $V$ be a strongly rational holomorphic vertex operator superalgebra and $G$ a finite solvable group
acting faithfully on $V$ by even automorphisms that fix the conformal vector, with
$G\cap\langle\sigma\rangle=1$. If $\alpha(V,G)=0$, then every trivialisation determines a strongly
rational holomorphic vertex operator superalgebra $W=V\orb G\supseteq V^G$ with the same conformal
vector, such that, as $\Z/2$- and $L_0$-graded spaces,
\[
  W=\bigoplus_{[g]}V(g)^{C(g)},\qquad W_R=\bigoplus_{[g]}V(g\sigma)^{C(g)},
\]
the sums running over the conjugacy classes of $G$. Here, for each representative $g$, the centraliser
$C(g)$ acts on $V(g)$ and on $V(g\sigma)$ by genuine representations $\rho_g$ through $x$-intertwiners
($x\in C(g)$), and $\rho_1$ is the given action.
\end{hypothesis}

In Theorem~\ref{thm:index} only the following parts are used: for $g$ of even order, that the summand of
$W_R$ belonging to $[g]$ is a graded subspace of $V(g\sigma)$; for $g$ of odd order, the formula
$V(g\sigma)^{C(g)}$ with a genuine action. In Proposition~\ref{prop:Sdefect} the NS decomposition is used
for a hypothetical orbifold. In both places the parities of the summands play no role.

\subsection{Further standard inputs}\label{sec:standard}
\begin{itemize}
\item[(A8)] Lifts of lattice automorphisms \cite{Lep}, \cite[Ch.~10]{FLM}, \cite[Prop.~4.1]{BK}: the
  automorphisms of $V_\Lambda$ preserving $M(1)$ and inducing $g$ on $\Lambda\otimes\C$ are the maps
  $u\otimes e^\lambda\mapsto\eta(\lambda)\,gu\otimes e^{g\lambda}$, and $\widehat\Lambda[2].O(\Lambda)$
  consists of those with $\eta(\lambda)=\pm1$.
\item[(A9)] The NS character of a strongly rational holomorphic vertex operator superalgebra is
  $S$-invariant: $\ch W(-1/\tau)=\ch W(\tau)$. Indeed it spans a one-dimensional space of trace functions
  stable under $S$ \cite[Thms.~4.12, 7.7]{DZ}, and $\ch W(i)>0$.
\item[(A10)] Twisted modules of lattice vertex superalgebras \cite[Thm.~4.2, Rem.~4.4]{BK}: existence,
  structure and characters (Proposition~\ref{prop:fpf}); on a $g$-twisted module $e^{2\pi iL_0}$ times the
  parity implements $g^{\pm1}$; and modular covariance of twisted--twined traces for commuting pairs that
  generate a cyclic group \cite[Thms.~4.12, 7.7]{DZ}.
\end{itemize}

\section{Cyclic groups}\label{sec:cyclic}

This section is \cite[Ex.~4.4]{GJF} and the anomaly half of the proof of \cite[Thm.~5.6]{GJF}, with two
misprints corrected and one citation replaced by a proof (Remark~\ref{rem:GJFtypos}).

\begin{proposition}\label{prop:maj}
For every $G\subseteq\Aut(C)$, the Majorana layer of $\alpha(V,G)$ is the restriction of the character
$\varepsilon$. In particular it vanishes for every group of Table~\ref{tab:main}.
\end{proposition}

\begin{proof}
The first claim is (A6) with Lemma~\ref{lem:majnumber}. For the second, $\varepsilon$ is a character, so
generators suffice: $N_-(-1)=n$; $N_-(J)=N_-(K)=n/2=6k$ by \eqref{eq:cycles}; elements of odd order have
$\varepsilon=1$; and $Jh$, $-h$ are products of these.
\end{proof}

\begin{proposition}\label{prop:frame}
Let $g\in\Aut(C)$ have order $m$, no fixed vector, and $\varepsilon(g)=1$. Then
$\alpha(V,\langle g\rangle)=a\,\nu_m$ with $a/m'\equiv L\rho_g\pmod1$.
\end{proposition}

\begin{proof}
$Z(1,1)=\Theta_\Lambda/\eta^n$ is invariant under $S$ and $T^L$ ($\Lambda$ is unimodular, $L$ is even,
and $Ln/24=Lk/2\in\Z$), so by (A7) $e_m(\alpha)$ is the eigenvalue of $ST^LS^{-1}$ on
$Z(1,g)=1/\eta_g$ (Proposition~\ref{prop:fpf}). Since $\eta(-1/\tau)=\sqrt{-i\tau}\,\eta(\tau)$ and
$\sum_kd_k=0$, $(1/\eta_g)|S=\prod_kk^{d_k/2}\prod_k\eta(\tau/k)^{-d_k}$. Every $k$ with $d_k\ne0$
divides $m$ and hence $L$, so $T^L$ multiplies this by
$\exp(-2\pi i\frac L{24}\sum_kd_k/k)=e^{2\pi iL\rho_g}e^{-2\pi iLn/24}=e^{2\pi iL\rho_g}$. Conclude with
Lemma~\ref{lem:calib}(2).
\end{proof}

\begin{proposition}\label{prop:odd}
Let $g\in\Aut(C)$ have odd order $m$. Then $\alpha(V,\langle g\rangle)=a\,\nu_m$ with
$a/m\equiv2m\rho_g\pmod1$. If $h\in\Aut(C)$ has order $3$ and no fixed coordinate, then
$\alpha(V,\langle h\rangle)=2k\,\nu_3$.
\end{proposition}

\begin{proof}
By (A8), $s(g)$ acts on $e^\lambda$ for $\lambda\in\Lambda^g$ by a sign $\eta(\lambda)$ with
$\eta(\lambda)^m=1$, so $\eta=1$ on $\Lambda^g$ ($s(g)$ is \emph{standard} \cite{Lep}), and
$Z(1,g)=\Theta_{\Lambda^g}/\eta_g$. Poisson summation gives
$\Theta_{\Lambda^g}(-1/\tau)=(-i\tau)^{r/2}\operatorname{vol}(\Lambda^g)^{-1}\Theta_{(\Lambda^g)^*}(\tau)$,
$r=\operatorname{rk}\Lambda^g$. As $\Lambda^g=\Lambda\cap(\Lambda\otimes\R)^g$ is primitive and $\Lambda$
unimodular, $(\Lambda^g)^*=\pi_g(\Lambda)$ for the orthogonal projection
$\pi_g=\frac1m\sum_jg^j$; and $m|\pi_g\lambda|^2=\sum_j\langle\lambda,g^j\lambda\rangle\in\Z$, so
$T^{2m}$ fixes $\Theta_{(\Lambda^g)^*}$. The eta factor behaves as in Proposition~\ref{prop:frame}, the
powers of $-i\tau$ cancel, and $Z(1,1)$ is invariant under $ST^{2m}S^{-1}$. This gives $a/m\equiv2m\rho_g$.
For $h$ the cycles have length $3$ and $\sigma=1$ (as $h^3=1$), so the Frame shape is $3^{n/3}$,
$\rho_h=n/27$, and $2m\rho_h=6n/27=8k/3$; hence $a\equiv8k\equiv2k\pmod3$. The Majorana layer vanishes
since $m$ is odd.
\end{proof}

\begin{corollary}\label{cor:cyclic}
For the elements used in Table~\ref{tab:main} the anomaly of the cyclic group they generate is $2k\,\nu_m$
in $\rSH(\Z/m)\cong\Z/m'$; it vanishes iff $m\mid k$:
\[
\begin{array}{lccccc}
\toprule
\text{element} & m & \text{signed cycles} & \rho_g & m' & \alpha\\
\midrule
-1 & 2 & 1\text{-cycles},\ \sigma=-1 & 3k/4 & 4 & 2k\\
x\text{ with }x^2=-1\ (J,K,JK) & 4 & 2\text{-cycles},\ \sigma=-1 & 9k/16 & 8 & 2k\\
-h & 6 & 3\text{-cycles},\ \sigma=-1 & 19k/36 & 12 & 2k\\
Jh & 12 & 6\text{-cycles},\ \sigma=-1 & 73k/144 & 24 & 2k\\
h & 3 & 3\text{-cycles},\ \sigma=1 & 4k/9 & 3 & 2k\\
\bottomrule
\end{array}
\]
\end{corollary}

\begin{proof}
The cycle types follow from \eqref{eq:cycles}: $(-h)^3=-1$, where $1$-cycles are excluded because
$(-h)^2=h^2$ has no fixed coordinate; and $(Jh)^6=-1$ with $(Jh)^4=h$, which excludes $2$-cycles of
$Jh$ in the same way. For the even types all cycles have length
$\ell=m/2$ and $\sigma=-1$, so $\rho_g=\frac n{24}(1+\frac1{2\ell^2})=\frac k2+\frac k{m^2}$, and
Proposition~\ref{prop:frame} gives $a\equiv m^2k+2k\equiv2k\pmod{2m}$, as $mk$ is even. The last row is
Proposition~\ref{prop:odd}.
\end{proof}

\begin{remark}[comparison with \cite{GJF}]\label{rem:GJFtypos}
(1) For $m$ odd the exponent in \cite[Ex.~4.4]{GJF} should be $\frac{2\pi i}{24}\sum_kd_k\frac{2m}k(k^2-1)$
(the operator is $ST^{2m}S^{-1}$); for $m=3$ the conclusion of \cite{GJF} is unaffected. (2) In the proof
of \cite[Thm.~5.6]{GJF} each block contributes $6$ times its number of sign changes to
$\langle\chi,(1-g)\chi\rangle$ (Lemma~\ref{lem:majnumber}), not $3(m/2)-2\times\mathrm{odd}$; the
conclusion $\langle\chi,(1-g)\chi\rangle\in4\Z$ is unaffected. (3) For $m=3$, \cite{GJF} cite
\cite[Ch.~5]{Moller} for the Frame-shape formula of a standard lift; that source treats even lattices,
while $\Lambda(C)$ is odd. Proposition~\ref{prop:odd} covers every unimodular lattice. (4) The extra
hypothesis of \cite[Ex.~4.4]{GJF} that the lift has order $m$ is automatic here, since all lifts of a
$g$ without fixed vector are conjugate to $s(g)$.
\end{remark}

\section{The quaternion group}\label{sec:Q8}

Throughout this section $C$, $J$, $K$ are as in Theorem~\ref{thm:Q8}, $n=12k$, and
$\alpha_C:=\alpha(V_{\Lambda(C)},Q_8)\in\SH(Q_8)$. The six elements $\pm J,\pm K,\pm JK$ satisfy $x^2=-1$
(note $(JK)^2=-J^2K^2=-1$), so none of them has a fixed coordinate or a fixed vector, and
$Q_8/\langle-1\rangle$ acts freely on the coordinates, with $3k$ orbits of size $4$.

\subsection{Restricted supercohomology of \texorpdfstring{$Q_8$}{Q8}}
We use the following standard facts \cite[XII.7]{CE}, \cite[IV]{AM}, \cite[VI.9]{Brown}.

\begin{lemma}\label{lem:cohQ8}
\begin{enumerate}
\item $\Hc^2(Q_8;\Z/2)\cong(\Z/2)^2$, and $\Hc^3(Q_8;U(1))\cong\Hc^4(Q_8;\Z)\cong\Z/8$, generated by
  the second Chern class $e$ of the $2$-dimensional irreducible representation $\rho_2$.
\item For $u\in\{i,j,k\}$ the restriction $\Hc^3(Q_8;U(1))\to\Hc^3(\langle u\rangle;U(1))\cong\Z/4$ is
  onto. Under the identifications $u\mapsto1$, the three restrictions of $e$ are equal, because $e$ is
  invariant under the automorphism $i\mapsto j\mapsto k\mapsto i$ ($\rho_2$ is the only $2$-dimensional
  irreducible representation).
\end{enumerate}
\end{lemma}

Identify each $\rSH(\langle u\rangle)$, $u\in\{i,j,k\}$, with $\Z/8$ through $\nu_4$ and $u\mapsto1$, and
let $R=(\res_{\langle i\rangle},\res_{\langle j\rangle},\res_{\langle k\rangle})\colon\rSH(Q_8)\to
(\Z/8)^3$.

\begin{proposition}\label{prop:kernel}
$\ker R=\{0,\kappa\}$ with $\kappa:=\iota_1(4e)\ne0$. The class $\kappa$ is fixed by every automorphism of
$Q_8$.
\end{proposition}

\begin{proof}
$|\rSH(Q_8)|=32$ (Lemma~\ref{lem:calib}). In $\rSH(\langle u\rangle)\cong\Z/8$ the image of $\iota_1$ is
the subgroup $2\Z/8$ of order $4$, and by Lemma~\ref{lem:cohQ8} $R(\iota_1e)=(2c,2c,2c)$ for one odd $c$.
For a character $\chi\colon Q_8\to\{\pm1\}$ let $b_\chi$ be the anomaly of $Q_8$ acting on $\Fer(2)$
through $\chi$, i.e.\ by the rotation through $\pi$ where $\chi=-1$. It lies in $\rSH(Q_8)$ by (A6), and
by Lemma~\ref{lem:calib}(2) its restriction to $\langle u\rangle$ is $0$ if $\chi(u)=1$ and
$\nu_{4,2}=4\nu_4$ if $\chi(u)=-1$. So $R(b_{\chi_i})=(0,4,4)$, $R(b_{\chi_j})=(4,0,4)$,
$R(b_{\chi_k})=(4,4,0)$, where $\chi_u$ has kernel $\langle u\rangle$. In $(2\Z/8)^3\cong(\Z/4)^3$ the
element $(1,1,1)$ generates a subgroup of order $4$, $(0,2,2)$ is not in it, and $(2,0,2)$ is not of the
form $(t,t+2s,t+2s)$; so these four elements generate a subgroup of order $16$, and $|\ker R|\le2$. Finally
$\iota_1$ is injective, $4e\ne0$, and $R(\kappa)=4R(\iota_1e)=0$. The last claim holds because every
automorphism of $Q_8$ fixes $e$.
\end{proof}

\begin{lemma}\label{lem:layers}
$\alpha_C\in\rSH(Q_8)$; $\res_{\langle x\rangle}\alpha_C=2k\,\nu_4$ for $x\in\{\pm J,\pm K,\pm JK\}$, and
$\res_{\langle-1\rangle}\alpha_C=2k\,\nu_2$. Hence $\alpha_C\ne0$ if $4\nmid k$, and $\alpha_C\in\{0,\kappa\}$
if $4\mid k$.
\end{lemma}

\begin{proof}
Proposition~\ref{prop:maj}, Corollary~\ref{cor:cyclic}, (A2) and (A5), and
Proposition~\ref{prop:kernel}.
\end{proof}

\subsection{Twisted sectors}
Let $\theta=s(-1)$; on each frame factor $V_{\sqrt3\Z}$ it is the reflection lift. We use the
following $q$-series, all carrying $q^{L_0-c/24}$:
\begin{gather*}
  e(\tau)=\frac{\eta(\tau)^n}{\eta(2\tau)^n},\quad
  f(\tau)=\frac{\eta(2\tau)^{n/2}}{\eta(4\tau)^{n/2}},\quad
  E(\tau)=2^{n/2}q^{n/48}\prod_{j\text{ odd}}(1-q^{j/2})^{-n},\\
  g(\tau)=2^{n/4}q^{n/192}\prod_{j\text{ odd}}(1-q^{j/4})^{-n/2},\quad
  F(\tau)=q^{n/48}\prod_{r\text{ odd}}(1+q^r)^{-n/2}.
\end{gather*}
By Proposition~\ref{prop:fpf}, $e=Z(1,-1)$ and $f=Z(1,x)$ for every $x$ with $x^2=-1$ (Frame shapes
$1^{-n}2^n$ and $2^{-n/2}4^{n/2}$).

\begin{lemma}\label{lem:minus}
\begin{enumerate}
\item $V(-1)=S[\theta]\otimes T$, where $S[\theta]$ is the Fock space of the Heisenberg modes $\alpha(r)$,
  $r\in\frac12+\Z$, and $T$ is the spinor module of the Clifford algebra $\mathrm{Cl}(n)$ generated by odd
  zero modes $\gamma_1,\dots,\gamma_n$, one for each frame factor; $\dim T=2^{n/2}$, all of $T$ has weight
  $n/16$, and $\ch V(-1)=E$. $V(-1)$ is irreducible already as a module over the frame.
\item For $g\in Q_8$, every operator on $V(-1)$ implementing $g$ is a multiple of $S[g]\otimes\phi_g$,
  where $S[g]$ is the action of $g$ on the Fock space and $\phi_g\in\mathrm{Pin}(n)$ satisfies
  $\phi_g\gamma_i\phi_g^{-1}=s_i\gamma_{p(i)}$.
\item There is a representation $\tilde\rho\colon Q_8\to GL(T)$ with $\tilde\rho(g)\in\C^\times\phi_g$ and
  $T\cong(\rho_2\oplus2\cdot\mathbf 1)^{\otimes3k}$. Hence $t_x:=\Tr_T\tilde\rho(x)=2^{3k}$ for
  $x\in\{\pm J,\pm K,\pm JK\}$ and $\Tr_T\tilde\rho(-1)=0$. Every representation $\rho$ with
  $\rho(g)\in\C^\times\phi_g$ is $\psi\otimes\tilde\rho$ for a character $\psi\colon Q_8\to\{\pm1\}$.
\item For $x\in\{\pm J,\pm K,\pm JK\}$, $V(x)=S[x]\otimes T_x$ with $T_x$ the spinor module of
  $\mathrm{Cl}(n/2)$, one zero mode for each signed $2$-cycle of $x$; $\dim T_x=2^{n/4}$, the
  superdimension of $T_x$ is $0$, the weights lie in $h_x+\frac14\Z_{\ge0}$ with $h_x=3n/64$, and
  $\ch V(x)=g$.
\item For $x\in\{\pm J,\pm K,\pm JK\}$ the trace of $S[x]$ on $S[\theta]$, weighted by
  $q^{L_0-c/24}$ with $L_0$ including the vacuum energy $n/16$ of $V(-1)$, is $F$.
\end{enumerate}
\end{lemma}

\begin{proof}
(1) On the frame, $\sqrt3\Z^n$ has a unique irreducible $\theta$-twisted module, $S[\theta]\otimes T$ with
$\dim T=|\sqrt3\Z^n/2\sqrt3\Z^n|^{1/2}=2^{n/2}$, since $n$ is even \cite[Thm.~4.2, Rem.~4.4]{BK}; its
zero modes are generated by the $\gamma_i$, which anticommute. $\theta$-twisted frame modules are
semisimple \cite[Thm.~4.2]{BK}, and $V(-1)$ has a ground-state space of dimension
$|\det(1-(-1))|^{1/2}=2^{n/2}$ (Proposition~\ref{prop:fpf}, $\varepsilon(-1)=1$), so as a frame module it
is this irreducible module. The vacuum energy is $\rho_{-1}=n/16$. (2) An implementing operator
conjugates the Heisenberg modes by $g$ and the zero modes as stated; by Schur's lemma for the
irreducible frame module it is unique up to a scalar. (3) On a $Q_8$-orbit $O$ of coordinates, $Q_8$ acts
on $\R^4_O$ by signed permutations, permuting $\pm e_p$ simply transitively; sending $x\in Q_8$ to
$xe_p$ for a fixed $p\in O$ identifies $\R^4_O$ with the quaternions, $Q_8$ acting by left
multiplication. So $Q_8$ acts through $SU(2)_L\subset SO(4)$ (or $SU(2)_R$), which lifts isomorphically
to $\mathrm{Spin}(4)=SU(2)_L\times SU(2)_R$; on the spinor module $S^+\oplus S^-=\C^2\oplus\C^2$ of
$\mathrm{Cl}(4)$ the lift acts as $\rho_2\oplus2\cdot\mathbf 1$. The lifts are even Clifford elements,
so they act factorwise on $T=\bigotimes_OS_O$ without Koszul signs, and their tensor product is a
representation $\tilde\rho$. The character of $\rho_2$ is $0$ on elements of order $4$ and $-2$ at $-1$,
so $t_x=2^{3k}$ and $\Tr\tilde\rho(-1)=0$. Two representations that agree up to scalars differ by a
character of $Q_8$, which is $\pm1$-valued. (4) As in (1), with $\sqrt3\Z^n/(1-x)\sqrt3\Z^n\cong
(\Z/2)^{n/2}$ generated by one $\sqrt3e_i$ for each $2$-cycle, and $|\det(1-x)|=2^{n/2}$; the vacuum energy
is $\rho_x=3n/64$ and the Fock modes lie in $\pm\frac14+\Z$. (5) $x$ has the eigenvalues $\pm i$ on
$\Lambda\otimes\C$, each $n/2$ times, so the Fock trace is $\prod_{r\in\frac12+\Z_{\ge0}}(1+q^{2r})^{-n/2}$,
and $q^{n/16-n/24}=q^{n/48}$.
\end{proof}

\begin{lemma}\label{lem:eta}
Write $(\phi|S)(\tau)=\phi(-1/\tau)$ and $(\phi|T)(\tau)=\phi(\tau+1)$. Then
\[
  e|S=E,\qquad f|S=g,\qquad (g|T^2)|S=N\cdot F,\quad\text{where }
  N=\bigl(\sqrt2\,e^{2\pi i/48}\bigr)^{n/2}=2^{3k}e^{2\pi ik/8}.
\]
Equivalently $f(\tau/(1-2\tau))=N\cdot F(\tau)$.
\end{lemma}

\begin{proof}
The first two follow from $\eta(-1/\tau)=\sqrt{-i\tau}\,\eta(\tau)$; for example
\[
  f(-1/\tau)=2^{n/4}\bigl(\eta(\tau/2)/\eta(\tau/4)\bigr)^{n/2}=g(\tau).
\]
For the third, use
$\eta(z+1)=e^{2\pi i/24}\eta(z)$ and Euler's identity
\[
  \eta\bigl(w+\tfrac12\bigr)=e^{2\pi i/48}\,\frac{\eta(2w)^3}{\eta(w)\,\eta(4w)},\qquad\text{i.e.}\qquad
  \prod_m\bigl(1-(-q)^m\bigr)=\frac{\prod_m(1-q^{2m})^3}{\prod_m(1-q^m)\prod_m(1-q^{4m})}.
\]
With $w=\tau/4$ they give
\[
  \frac{\eta(\tau/2+1)}{\eta(\tau/4+\frac12)}=e^{2\pi i/48}\,\frac{\eta(\tau)\eta(\tau/4)}{\eta(\tau/2)^2},
  \qquad\text{so}\qquad
  g(\tau+2)=N\Bigl(\frac{\eta(\tau)\eta(\tau/4)}{\eta(\tau/2)^2}\Bigr)^{n/2}=N\,F(-1/\tau),
\]
using $F=\bigl(\eta(\tau)\eta(4\tau)/\eta(2\tau)^2\bigr)^{n/2}$ and the $S$-law of $\eta$ (the factors
$\sqrt{-i\tau}$ cancel). Applying $S$ gives $(g|T^2)|S=N\,F$, and $f|ST^2S=(g|T^2)|S$.
\end{proof}

\subsection{The character of a would-be orbifold}

\begin{proposition}\label{prop:Sdefect}
Assume (A1), (A9), (A10) and Hypothesis~ORB. If $4\mid k$ and the $Q_8$-action on $V_{\Lambda(C)}$ is
non-anomalous, then $8\mid k$.
\end{proposition}

\begin{proof}
Let $W=V\orb Q_8$ for some trivialisation. The conjugacy classes of $Q_8$ are $\{1\}$, $\{-1\}$,
$\{\pm J\}$, $\{\pm K\}$, $\{\pm JK\}$, with $C(\pm1)=Q_8$ and $C(x)=\langle x\rangle$ otherwise. We
compute the contribution of each class to $\ch W$.

\emph{Untwisted sector.} $\frac18\sum_{g\in Q_8}Z(1,g)=\frac18\bigl(Z(1,1)+e+6f\bigr)$.

\emph{The class of $-1$.} By Lemma~\ref{lem:minus}(2),(3) and genuineness,
$\rho_{-1}(g)=\psi(g)\,S[g]\otimes\tilde\rho(g)$ for a character $\psi$. The trace of $\rho_{-1}(g)$ is
$E$ for $g=1$, $0$ for $g=-1$, and $\psi(x)t_xF$ for $g=\pm x$ (note $\psi(-x)=\psi(x)$, since $-1$ is a
commutator, and $t_{-x}=t_x$). So the contribution is $\frac18E+\frac14T_\psi F$ with
$T_\psi=2^{3k}\bigl(\psi(J)+\psi(K)+\psi(JK)\bigr)$.

\emph{The class of $x\in\{J,K,JK\}$.} By (A10) and Schur's lemma, $\rho_x(x)=\epsilon\,e^{2\pi iL_0}P$ for a
scalar $\epsilon$, $P$ the parity. On the weight space of weight $h_x+j/4$, which is (Fock level $j/4$)
$\otimes T_x$, odd powers of $\rho_x(x)$ have trace $0$ because $T_x$ has superdimension $0$, and
$\rho_x(x)^2$ is the scalar $\epsilon^2e^{4\pi ih_x}(-1)^j$. So the projector onto the invariants has
trace $\frac14(1+\epsilon^2e^{4\pi ih_x}(-1)^j)$ times the dimension, which is nonzero for $j=0$ or $j=1$;
both levels are nonzero. Hence $W_x:=V(x)^{\langle x\rangle}\ne0$. On $W_x\subseteq W$ the operator
$e^{2\pi iL_0}$ is $\pm P$ ($W$ is a vertex operator superalgebra) and $\rho_x(x)=1$, so $\epsilon=\pm1$
and $\rho_x(x^2)=e^{4\pi iL_0}$, whose trace is $e^{4\pi ic/24}g(\tau+2)=g(\tau+2)$. The contribution is
$\frac14\bigl(g(\tau)+g(\tau+2)\bigr)$.

\emph{$S$-invariance.} Altogether
\[
  \ch W=\tfrac18Z(1,1)+\tfrac18e+\tfrac34f+\tfrac18E+\tfrac14T_\psi F+\tfrac34\bigl(g+g|T^2\bigr).
\]
$Z(1,1)=\Theta_\Lambda/\eta^n$ is $S$-invariant, and by Lemma~\ref{lem:eta} $e\leftrightarrow E$,
$f\leftrightarrow g$, $g|T^2\mapsto NF$ and $F\mapsto(g|T^2)/N$ under $S$. Hence
\[
  \ch W|S-\ch W=\tfrac14\bigl(3N-T_\psi\bigr)\bigl(F-(g|T^2)/N\bigr).
\]
By (A9) the left side vanishes. The leading exponents of $F$ and $g|T^2$ are $n/48$ and $n/192$, so
$F\ne(g|T^2)/N$ and $T_\psi=3N$, i.e.\ $\psi(J)+\psi(K)+\psi(JK)=3e^{2\pi ik/8}$. The left side is $3$
($\psi=1$) or $-1$, so $\psi=1$ and $e^{2\pi ik/8}=1$.
\end{proof}

\subsection{Proof of Theorem~\ref{thm:Q8}}
If $4\nmid k$, $\res_{\langle J\rangle}\alpha_C=2k\nu_4\ne0$ (Lemma~\ref{lem:layers}). If $k\equiv4\pmod8$,
then $\alpha_C\in\{0,\kappa\}$ by Lemma~\ref{lem:layers}, and $\alpha_C=0$ is excluded by
Proposition~\ref{prop:Sdefect}; so $\alpha_C=\kappa\ne0$. Let $8\mid k$. Let $D$ be a self-dual code of
length $48$ with $J_D,K_D\in\Aut(D)$ satisfying the same relations; for instance the sum of twelve
$Q_8$-invariant tetracodes, $Q_8$ acting on each $4$-set by left multiplication on quaternions
\cite[Lemma~6.2]{I}. Then $Q_8$ acts on $C\oplus D$ through $(J,J_D)$ and $(K,K_D)$, and
$s_{C\oplus D}=s_C\otimes s_D$, so $\alpha_{C\oplus D}=\alpha_C+\alpha_D$ by (A4). Both $C\oplus D$ and
$D$ have $k\equiv4\pmod8$, so both anomalies equal $\kappa$, which does not depend on the identification
of the groups with $Q_8$ (Proposition~\ref{prop:kernel}). Hence $\alpha_C=0$. \qed

\medskip
The proof shows more: $\alpha_C=(k/4)\,\kappa$ when $4\mid k$.

\begin{remark}[free fermions]\label{rem:freefermions}
For a $Q_8$-invariant sum $D_k$ of tetracodes of length $12k$ the lattice is $\Z^{12k}$ (the only odd
unimodular lattice of rank $4$ is $\Z^4$), so $V_{\Lambda(D_k)}=\Fer(8)^{\otimes3k}$. The weight-$\frac12$
space of each factor is spanned by the $e^{t/\sqrt3}$, $t$ a nonzero tetracode word, which $Q_8$ permutes
up to sign and freely; so it is the real regular representation
$\Reg=\mathbf 1\oplus\chi_i\oplus\chi_j\oplus\chi_k\oplus\mathbb H$. Let $a=\alpha(\Fer(\Reg))$. Then
$\alpha_{D_k}=3k\,a$, so $12a=\kappa$ and $24a=0$ by the proof above; as $|\rSH(Q_8)|=32$, $a$ has order
exactly $8$ and $4a=\kappa$. Moreover $\alpha_C=3k\,a$ for every $C$: the difference $\delta(C)=\alpha_C-
3k\,a$ lies in $\ker R$ ($\Reg$ restricted to $\langle i\rangle$ rotates two planes through $\pi/2$ and one
through $\pi$, so its anomaly there is $(1+1+4)\nu_4$, and $3k\cdot6\equiv2k\pmod8$), it is additive, it
vanishes on the $D_k$, and it vanishes for $4\mid k$; adding a $D_j$ with $j\equiv-k\pmod4$ gives
$\delta(C)=0$ in general.

In spin-bordism terms the anomaly of fermions in a real representation $\rho$ on a spherical space form
$S^3/\Gamma$ with spin structure $\epsilon\colon\Gamma\to\{\pm1\}$ is
$\frac12(\eta(\rho)-\dim\rho\,\eta(\mathbf 1))\bmod1$, with Donnelly's
$\eta(\rho)=-\frac1{|\Gamma|}\sum_{g\ne1}\epsilon(g)\chi_\rho(g)/(2-\operatorname{tr}g)$ \cite{Donnelly}. For
$\rho=\Reg$ this is $\frac58$ on $S^3/Q_8$ for all four spin structures, $\frac14$ on $S^3/\langle
i\rangle$ and $\frac12$ on $\mathbb{RP}^3$; one Majorana fermion odd under $\Z/2$ gives $\pm\frac18$ on
$\mathbb{RP}^3$. This is consistent with $a$ of order $8$ and with $4a$ vanishing on all cyclic
subgroups. We do not use these values.

The twisted sector of Lemma~\ref{lem:minus} also gives a direct invariant. Normalise the lifts of $J$, $K$,
$JK$ on $V(-1)$ by transporting their twining characters on $V$ along $ST^2S$ (Lemma~\ref{lem:eta}); for
$4\mid k$ these are the canonical lifts of the non-anomalous cyclic subgroups. They differ from
$\tilde\rho$ by scalars, and $\hat J\hat K=\mu\,\widehat{JK}$ for a scalar $\mu$; with the conventions of
the computations in \cite{repo} one finds $\mu=e^{-2\pi ik/8}=\exp(2\pi i\cdot3k\cdot\frac58)$, the value of
$3k\,a$ on $S^3/Q_8$. Finally, $\alpha_C=3k\,a$ can
also be derived by gauging: a sum $C_0$ of $Q_8$-invariant tetracodes is a multi-neighbour of $C$, and
$V_{\Lambda(C_0)}=V_{\Lambda(C)}\orb\widehat W$ for a $3$-group $\widehat W$ normalised by $Q_8$, so a
coprime-gauging lemma \cite{Naidu,Uribe,Tachikawa} transports the anomaly from $C$ to $C_0$. That route
needs the super version of the gauging lemma, which the proof above avoids.
\end{remark}

\begin{remark}
For $k\equiv4\pmod8$ a non-anomalous $Q_8$ would produce an orbifold of index $3$, which is not divisible
by $24/\gcd(k,24)\in\{6,2\}$. So Theorem~\ref{thm:Q8} is what the Stolz--Teichner conjecture predicts.
\end{remark}

\section{The group \texorpdfstring{$Q_8\times\Z/3$}{Q8 x Z/3}}\label{sec:Q8Z3}

\begin{lemma}[{\cite[III.9--III.10]{Brown}}]\label{lem:sylow}
Let $G$ be finite, $M$ a trivial $G$-module, $q\ge1$, and $P$ a Sylow $p$-subgroup. Then $\res^G_P$ is
injective on the $p$-primary part of $\Hc^q(G;M)$, and on all of $\Hc^q(G;\Z/2)$ if $p=2$. Hence
$\beta\in\Hc^q(G;M)$ vanishes iff all its Sylow restrictions vanish.
\end{lemma}

\begin{proposition}\label{prop:sylow}
Let $G$ be finite and $\alpha\in\SH(G)$. Then $\alpha=0$ iff $\res^G_P\alpha=0$ for every Sylow subgroup
$P$.
\end{proposition}

\begin{proof}
Suppose all Sylow restrictions vanish, and let $P_2$ be a Sylow $2$-subgroup. By naturality
$\res_{P_2}\maj(\alpha)=\maj(\res_{P_2}\alpha)=0$, so $\maj(\alpha)=0$ by Lemma~\ref{lem:sylow}, and
$\alpha=\iota_2(\alpha')$ with $\alpha'\in\rSH(G)$; since $\iota_2$ is injective for every subgroup, all
Sylow restrictions of $\alpha'$ vanish. In the same way $\gw(\alpha')=0$, so $\alpha'=\iota_1(\beta)$ with
$\beta\in\Hc^3(G;U(1))$ whose Sylow restrictions vanish, and $\beta=0$ by Lemma~\ref{lem:sylow}.
\end{proof}

\begin{proof}[Proof of Theorem~\ref{thm:Q8Z3}]
The Sylow subgroups of $G=Q_8\times\langle h\rangle$ are $Q_8$ and $\langle h\rangle$, both normal. The
restriction to $Q_8$ is $\alpha_C$, which vanishes iff $8\mid k$ (Theorem~\ref{thm:Q8}). An element of
order $3$ without fixed coordinate consists of $3$-cycles, so the restriction to $\langle h\rangle$ is
$2k\,\nu_3$ (Proposition~\ref{prop:odd}), which vanishes iff $3\mid k$. Conclude with
Proposition~\ref{prop:sylow}.
\end{proof}

The same argument shows that $(\res_{Q_8},\res_{\Z/3})\colon\SH(Q_8\times\Z/3)\to\SH(Q_8)\oplus\SH(\Z/3)$
is an isomorphism (inflations give the inverse): there are no mixed terms.

\section{The orbifold index}\label{sec:index}

Throughout, $C$ and $G$ are as in Theorem~\ref{thm:index}, $V=V_{\Lambda(C)}$, $G$ acts non-anomalously,
and $W=V\orb G$ is an orbifold as in Hypothesis~ORB. Since the anomaly vanishes, so does its Majorana
layer, hence $\varepsilon=1$ on $G$ (Proposition~\ref{prop:maj}).

\begin{lemma}\label{lem:groups}
\begin{enumerate}
\item Every commuting pair in $G$ generates a cyclic subgroup.
\item If $g\in G$ has even order $m$, then $g^{m/2}=-1$.
\item Every $g\in G\setminus\{\pm1\}$ has no fixed coordinate; so $\operatorname{tr}_{\R^n}g=0$.
\end{enumerate}
\end{lemma}

\begin{proof}
(1) is (i). (2) $\langle g\rangle$ contains exactly one involution, $g^{m/2}$, and it is $-1$ by (ii).
(3) A coordinate fixed by $g$ is fixed by every power of $g$. By (iii) the order of $g$ is $2^a3^b$ with
$b\le1$. If $b=1$, then $g^{2^a}$ is a nontrivial element of $\langle h\rangle$, which has no fixed
coordinate. If $b=0$, then $a\ge2$, since $g\ne\pm1$ and $-1$ is the only involution; so
$g^{2^{a-1}}=-1$ by (2), and a coordinate fixed by $g$, with $ge_i=\pm e_i$, would give
$-e_i=g^{2^{a-1}}e_i=e_i$. A signed permutation matrix without fixed coordinate has trace $0$.
\end{proof}

\begin{lemma}\label{lem:vacuum}
Let $g\in G$ have even order $m$. Then the eigenvalues of $g$ on $\R^n$ are $e^{2\pi is/m}$ with $s$ odd,
each with multiplicity $2n/m$; $g$ has no fixed vector; and the lowest weight of $V(g\sigma)$ is
\[
  E_g=\rho_g=\frac n{24}+\frac n{12m^2}>\frac c{24}.
\]
\end{lemma}

\begin{proof}
By Lemma~\ref{lem:groups} the character of $\langle g\rangle$ on $\R^n$ is $n$ at $1$, $-n$ at
$g^{m/2}=-1$, and $0$ elsewhere, so $e^{2\pi is/m}$ has multiplicity $\frac nm(1-(-1)^s)$. With $m=2p$,
$\sum_{s\text{ odd}<m}\frac sm(1-\frac sm)=\frac{2p^2+1}{12p}$, so $\rho_g=\frac n4\cdot\frac
2m\cdot\frac{2p^2+1}{12p}=\frac n{24}+\frac n{12m^2}$. By Proposition~\ref{prop:fpf} ($\varepsilon(g)=1$)
the lowest weight of $V(g\sigma)\cong V(g)$ is $\rho_g$.
\end{proof}

\begin{lemma}\label{lem:evenness}
Every $V(g)$ and $V(g\sigma)$, $g\in G$, is of $\C$-type. If $g$ has odd order and $x\in C(g)$, every
$x$-intertwiner on $V(g\sigma)$ is even.
\end{lemma}

\begin{proof}
The Majorana layer, which records which $V(g)$ are of Clifford type, is $\varepsilon=1$; and $V_R$ is of
$\C$-type, with $V(g\sigma)\cong V(g)\boxtimes V_R$ in the $G\times\langle\sigma\rangle$-crossed
supercategory of twisted modules \cite[Thms.~4.15, 4.17]{McRae}, in which each component has one simple
object, invertible. For a $\C$-type simple object $M$ let $p(M,x)\in\Z/2$ be the parity of an
$x$-intertwiner; it is well defined, $p(V,x)=0$, and $p$ is additive under $\boxtimes$ (an isomorphism
$M\boxtimes N\cong P$ of either parity enters twice). So $a\mapsto p(V(g^a),x)$ is a homomorphism from
$\Z/\operatorname{ord}(g)$ to $\Z/2$, trivial for odd order, and $p(V(g\sigma),x)=p(V_R,x)$. Finally
$p(V_R,x)=0$: an $x$-intertwiner maps $e^{\nu_w}$ to a multiple of $e^{\nu_{xw}}$, whose parity is
$f(xw)=\varepsilon(x)f(w)=f(w)$ (if $C$ has no maximal codeword, $\Index(C)=0$ and there is nothing to
prove).
\end{proof}

\begin{lemma}\label{lem:untwisted}
Let $A$ be an $x$-intertwiner on $V_R$, $x\in G$, and $Z_{RR}(1,A)=\str_{V_R}Aq^{L_0-c/24}$.
\begin{enumerate}
\item $Z_{RR}(1,A)=0$ if $x$ has even order.
\item If $x$ has order $3$ and $A^3=1$, then $Z_{RR}(1,A)=\zeta\,T_x(C)$ for some $\zeta^3=1$.
\end{enumerate}
\end{lemma}

\begin{proof}
By Proposition~\ref{prop:VC}(5) only the ground states $e^{\nu_w}$ contribute, and $A$ maps
$\C e^{\nu_w}$ onto $\C e^{\nu_{xw}}$, since it conjugates the Heisenberg zero modes by $x$. (1) A
maximal codeword fixed by $x$ would be fixed by $x^{m/2}=-1$, but $-w\ne w$. (2) The lift $s(x)$ acts on
$e^\alpha$, $\alpha\in\Lambda^x$, by a sign whose cube is $1$, hence trivially. Any two $x$-fixed ground
states $e^{\nu_w},e^{\nu_{w'}}$ have momenta differing by some $\alpha\in\Lambda^x$, and the leading term
of $Y_R(e^\alpha,z)e^{\nu_w}$ is a nonzero multiple of $e^{\nu_w+\alpha}$; since
$AY_R(e^\alpha,z)A^{-1}=Y_R(e^\alpha,z)$, $A$ has one eigenvalue $\zeta$ on all $x$-fixed ground states,
and $\zeta^3=1$. The parities are $f(w)$, so $Z_{RR}(1,A)=\zeta\sum_{xw=w}f(w)=\zeta\,T_x(C)$.
\end{proof}

\begin{proof}[Proof of Theorem~\ref{thm:index}]
$W$ contains $V^G\ni\tau$ with the same conformal vector, so it is $\cN=1$ with the supercurrent $\tau$,
and $G_0^2=L_0-c/24$ on $W_R$. Hence $\Index(W)=\str_{W_R}q^{L_0-c/24}$ is the superdimension of the
weight-$c/24$ part of $W_R$, which is the sum of the contributions of the classes $[g]$.

\emph{$g$ of even order.} The summand lies in $V(g\sigma)$, which has no states of weight $c/24$ by
Lemma~\ref{lem:vacuum}; it contributes $0$.

\emph{$g$ of odd order.} The summand is $V(g\sigma)^{C(g)}$, the image of the projector
\[
  e_g=\frac1{|C(g)|}\sum_{x\in C(g)}\rho_g(x),
\]
which is even by Lemma~\ref{lem:evenness}. Each $\rho_g(x)$ commutes with $G_0$,
since $x$ fixes $\tau$, so the contribution is $\frac1{|C(g)|}\sum_{x\in C(g)}Z_{RR}(g,x)$ with
$Z_{RR}(g,x)=\str_{V(g\sigma)}\rho_g(x)q^{L_0-c/24}$.

For $g=1$ this is $\frac1{|G|}\sum_xZ_{RR}(1,x)=\Index(C)/|G|$, by Lemma~\ref{lem:untwisted}: elements of
even order contribute $0$, and the elements of order $3$ are $h^{\pm1}$, with $T_{h^{\pm1}}=T_h=0$.

Let $g\in\langle h\rangle\setminus\{1\}$; then $C(g)=G$ by (iii). If $x\in G$ has even order, then
$\langle g,x\rangle$ is cyclic, and by modular covariance (A10) $Z_{RR}(g,x)$ is a multiple of the
$S$-transform of $Z_{RR}(x,g^{-1})$, which vanishes because $V(x\sigma)$ has no states of weight $c/24$.
If $x$ has odd order, then $x\in\langle g\rangle$; as $SL_2(\Z)\to SL_2(\Z/3)$ is onto and $SL_2(\Z/3)$
acts transitively on $(\Z/3)^2\setminus0$, some $\gamma\in SL_2(\Z)$ carries the pair $(g,x)$ to $(1,y)$
with $y\ne1$ in $\langle h\rangle$, and $Z_{RR}(1,y)=\zeta T_y(C)=0$. So all twisted classes contribute
$0$, and $\Index(W)=\pm\Index(C)/|G|$, the sign being the convention for the parity of $W_R$.
\end{proof}

\begin{remark}
For $G=Q_8$, and for cyclic $G$ of order $1$, $2$ or $4$, every class except $[1]$ has even order, so no
modular transformation of twisted--twined traces is used. Changing the trivialisation can change
individual twisted--twined terms by phases, but every term except $Z_{RR}(1,1)$ vanishes for every
normalisation, so the index is the same for all orbifolds.
\end{remark}

\section{Proof of the Main Theorem, and remarks}\label{sec:proof}

\begin{proof}[Proof of Theorem~\ref{thm:main}]
Let $m=\gcd(k,24)$. By \cite[Thms.~B, C]{I} there is a self-dual code $C$ of length $12k$ with
$\Index(C)=24$ and the automorphisms of Table~\ref{tab:main}: for $m\in\{1,2\}$ any code of
\cite[Thm.~B]{I}; for $3\mid k$ an $h$ of order $3$ without fixed coordinate and with $T_h(C)=0$
\cite[Thm.~C(a)]{I}; for $4\mid k$ elements $J,K$ with $J^2=K^2=-1$, $KJ=-JK$ \cite[Thm.~C(b)]{I}; and
for $12\mid k$ all three, commuting \cite[Thm.~C(c)]{I}. For $m\in\{1,2,3,4,6,12\}$ the group $G$ is cyclic,
generated by $1$, $-1$, $h$, $J$, $-h$, $Jh$, of order $m$; by Corollary~\ref{cor:cyclic} its anomaly is
$2k\nu_m=0$, as $m\mid k$. For $m=8$, $G=Q_8$ is non-anomalous by Theorem~\ref{thm:Q8}, and for $m=24$,
$G=Q_8\times\langle h\rangle$ by Theorem~\ref{thm:Q8Z3}. Each $G$ satisfies (i)--(iii) of
Theorem~\ref{thm:index}: its abelian subgroups are cyclic (for $Q_8\times\Z/3$ they are products of
subgroups of coprime orders); $-1=g^{m/2}$ resp.\ the centre of $Q_8$ is its only involution; and its
elements of odd order form $1$ or $\langle h\rangle$, central, with $T_h(C)=0$ ($T_{h^2}=T_h$, and
$(Jh)^4=h$). By Hypothesis~ORB the orbifold $W=V_{\Lambda(C)}\orb G$ exists, and by
Theorem~\ref{thm:index} it is a holomorphic $\cN=1$ vertex operator superalgebra of central charge
$12k$ and index $24/|G|=24/m$.
\end{proof}

\begin{remark}[reduction to $k\le24$]\label{rem:reduction}
If $W$ and $W'$ are holomorphic $\cN=1$ vertex operator superalgebras of central charges in $12\Z$, then
so is $W\otimes W'$, with supercurrent $\tau\otimes\mathbf 1+\mathbf 1\otimes\tau'$; its Ramond sector is
$W_R\otimes W'_R$, so $\Index(W\otimes W')=\pm\Index(W)\Index(W')$. With $W_{24}$ the orbifold of index $1$
at $k=24$ and $W_r$ that of index $24/\gcd(r,24)$ at $k=r\le24$, the theory $W_{24}^{\otimes q}\otimes W_r$
has central charge $12(24q+r)$ and index $24/\gcd(24q+r,24)$. So the Main Question needs only the codes
of length at most $288$; the chains of \cite{I} are used for finitely many steps.
\end{remark}

\begin{remark}[$k=24$ and TMF]\label{rem:TMF}
Let $W$ be the orbifold of index $1$ and central charge $288$. Under the Stolz--Teichner conjecture
\cite[Conj.~5.4]{GJF} it defines a class $[W]\in\pi_{576}\mathrm{TMF}$ mapping to $\pm\Delta^{24}$ in
$\mathrm{MF}_{288}$, so $[W]\mp\Delta^{24}$ lies in the kernel of $\pi_*\mathrm{TMF}\to\mathrm{MF}$, which
is the torsion. Torsion elements of $\pi_*\mathrm{TMF}$ are nilpotent by the nilpotence theorem \cite{DHS},
their images in $\mathrm{MU}_*\mathrm{TMF}$ being torsion in a torsion-free group. Hence
$[W]=\pm\Delta^{24}(1+t)$ with $t$ nilpotent, a unit of $\pi_*\mathrm{TMF}$; and a theory of central
charge $288$ can define a unit only if its index is $\pm1$. So at $k=24$ the Main Question asks for a
holomorphic $\cN=1$ theory representing the periodicity class $\Delta^{24}$ up to torsion, and
Theorem~\ref{thm:main} provides one.
\end{remark}

\begin{remark}[why the seed for $12\mid k$ has length $72$]
The codes for $12\mid k$ need $Q_8\times\Z/3$ and are built from a seed of length $72$
\cite[\S7.3]{I}. $Q_8$-invariant codes are $\F^2\otimes U$ with $U$ Lagrangian in $\F^{2r}$ for the
symplectic form $\sum_b(u_{2b}v_{2b+1}-u_{2b+1}v_{2b})$ (cf.\ \cite[App.~A.2]{I}), which makes their index
computable at length $72$. At length $36$, in all our samples, $Q_8\times\Z/3$-invariant codes with
$T_h=0$ have index divisible by $72$, which rules out seeds of that length for chains with index $24$;
at length $72$ this pattern fails, and a seed exists.
\end{remark}

\begin{remark}[formalisation]
The code-level content of Theorem~\ref{thm:main}, i.e.\ for every $k$ a self-dual code of length $12k$ and
index $24$ with the group data of Table~\ref{tab:main}, is formalised in Lean~4 with Mathlib in the
repository of \cite{I}, with the value $24$ of the length-$72$ seed sum as an explicit hypothesis. The
vertex-algebra side is not formalised.
\end{remark}

\appendix
\section{Computations}\label{app:checks}

The explicit identities of this paper were checked by computer, with the code data parsed from the
source of \cite{I}; the script and its log are in \cite{repo}.
\begin{itemize}
\item The Ramond ground states of Proposition~\ref{prop:VC}(4) (the per-coordinate minima, and the signed
  count $24$ on the Golay code), Lemma~\ref{lem:majnumber} and \eqref{eq:cycles} on random signed
  permutations, and the two expressions for $\rho_g$.
\item The calibration of Lemma~\ref{lem:calib} (theta quotients under $ST^LS^{-1}$, numerically); the
  eigenvalues of Propositions~\ref{prop:frame} and~\ref{prop:odd} numerically (eta quotients for
  $n=12,24$; the theta series of $\Lambda^h$ for an element $h$ of order $3$ without fixed coordinate on
  the Golay code and on $E_{36}$); and the anomalies $2k\nu_m$ of
  Corollary~\ref{cor:cyclic} exactly, for $k\le48$, and for every element of every group of
  Table~\ref{tab:main} on the codes of \cite{I}, together with $\varepsilon=1$ and the Frame shapes.
\item The cohomology of $Q_8$ in low degree over $\mathbb F_2$, and the order $16$ of the subgroup of
  Proposition~\ref{prop:kernel}.
\item Lemma~\ref{lem:minus}(3): explicit Pin lifts in $\mathrm{Cl}(4)$ on every $Q_8$-orbit of coordinates
  of the Golay code and of $E_{72}$ (even, unique up to scalars, with exactly four
  linearisations $\rho_2\oplus2\psi$), and on the whole spinor module of $\mathrm{Cl}(12)$ for the Golay
  code (traces $8\psi(x)$ at elements of order $4$ and $0$ at $-1$).
\item The three identities of Lemma~\ref{lem:eta} (numerically at several $\tau$, for $n=12,24,48$); the
  formula for $\ch W|S-\ch W$ in Proposition~\ref{prop:Sdefect} numerically for $k=4,8$ and all four
  $\psi$, and its zero set ($8\mid k$, $\psi=1$) exactly.
\item The eta invariants of Remark~\ref{rem:freefermions}, in exact rational arithmetic.
\item The eigenvalue multiplicities and the vacuum energies $E_g$ of Lemma~\ref{lem:vacuum} for every
  element of even order, and the group theory of Lemma~\ref{lem:groups}, for every group of
  Table~\ref{tab:main} on the codes of \cite{I}; and the arithmetic of Table~\ref{tab:main} and
  Remark~\ref{rem:reduction}.
\end{itemize}

\begin{thebibliography}{99}

\bibitem{AM}
A.~Adem and R.~J.~Milgram, \emph{Cohomology of Finite Groups}, 2nd ed., Grundlehren Math.\ Wiss.\
\textbf{309}, Springer, 2004.

\bibitem{BK}
B.~Bakalov and V.~G.~Kac, \emph{Twisted modules over lattice vertex algebras}, in: Lie Theory and Its
Applications in Physics V, World Sci., 2004, 3--26; arXiv:math/0402315.

\bibitem{Brown}
K.~S.~Brown, \emph{Cohomology of Groups}, Grad.\ Texts in Math.\ \textbf{87}, Springer, 1982.

\bibitem{CM}
S.~Carnahan and M.~Miyamoto, \emph{Regularity of fixed-point vertex operator subalgebras},
arXiv:1603.05645.

\bibitem{CE}
H.~Cartan and S.~Eilenberg, \emph{Homological Algebra}, Princeton Univ.\ Press, 1956.

\bibitem{DHS}
E.~S.~Devinatz, M.~J.~Hopkins and J.~H.~Smith, \emph{Nilpotence and stable homotopy theory I}, Ann.\ of
Math.\ \textbf{128} (1988), 207--241.

\bibitem{DL}
C.~Dong and J.~Lepowsky, \emph{The algebraic structure of relative twisted vertex operators}, J.\ Pure
Appl.\ Algebra \textbf{110} (1996), 259--295.

\bibitem{DNR}
C.~Dong, S.-H.~Ng and L.~Ren, \emph{Orbifolds and minimal modular extensions}, arXiv:2108.05225.

\bibitem{DRY}
C.~Dong, L.~Ren and M.~Yang, \emph{Super orbifold theory}, Adv.\ Math.\ \textbf{405} (2022), 108481;
arXiv:2104.03533.

\bibitem{DZ}
C.~Dong and Z.~Zhao, \emph{Modularity in orbifold theory for vertex operator superalgebras}, Comm.\
Math.\ Phys.\ \textbf{260} (2005), 227--256; arXiv:math/0411524.

\bibitem{Donnelly}
H.~Donnelly, \emph{Eta invariants for $G$-spaces}, Indiana Univ.\ Math.\ J.\ \textbf{27} (1978),
889--918.

\bibitem{FLM}
I.~Frenkel, J.~Lepowsky and A.~Meurman, \emph{Vertex Operator Algebras and the Monster}, Pure Appl.\
Math.\ \textbf{134}, Academic Press, 1988.

\bibitem{GJF17}
D.~Gaiotto and T.~Johnson-Freyd, \emph{Symmetry protected topological phases and generalized
cohomology}, J.\ High Energy Phys.\ \textbf{2019}, no.~5, 007; arXiv:1712.07950.

\bibitem{GJF}
D.~Gaiotto and T.~Johnson-Freyd, \emph{Holomorphic SCFTs with small index}, Canad.\ J.\ Math.\
\textbf{74} (2022), no.~2, 573--601; arXiv:1811.00589. (Numbering of statements refers to the arXiv
version.)

\bibitem{Hopkins}
M.~J.~Hopkins, \emph{Algebraic topology and modular forms}, Proc.\ ICM Beijing 2002, Vol.~I, 291--317.

\bibitem{I}
[Authors], \emph{The index of self-dual ternary codes}, companion paper (draft, 2026).

\bibitem{Kirillov}
A.~Kirillov~Jr., \emph{Modular categories and orbifold models}, Comm.\ Math.\ Phys.\ \textbf{229} (2002),
309--335.

\bibitem{Lep}
J.~Lepowsky, \emph{Calculus of twisted vertex operators}, Proc.\ Nat.\ Acad.\ Sci.\ USA \textbf{82}
(1985), 8295--8299.

\bibitem{Li}
H.~Li, \emph{Local systems of twisted vertex operators, vertex operator superalgebras and twisted
modules}, in: Moonshine, the Monster, and Related Topics, Contemp.\ Math.\ \textbf{193} (1996),
203--236; arXiv:q-alg/9504022.

\bibitem{McRae}
R.~McRae, \emph{Twisted modules and $G$-equivariantization in logarithmic conformal field theory},
Comm.\ Math.\ Phys.\ \textbf{383} (2021), 1939--2019; arXiv:1910.13226.

\bibitem{Moller}
S.~M\"oller, \emph{A cyclic orbifold theory for holomorphic vertex operator algebras and applications},
PhD thesis, TU Darmstadt, 2016; arXiv:1611.09843.

\bibitem{Naidu}
D.~Naidu, \emph{Categorical Morita equivalence for group-theoretical categories}, Comm.\ Algebra
\textbf{35} (2007), 3544--3565; arXiv:math/0605530.

\bibitem{repo}
[Repository URL to be added].

\bibitem{Tachikawa}
Y.~Tachikawa, \emph{On gauging finite subgroups}, SciPost Phys.\ \textbf{8} (2020), 015;
arXiv:1712.09542.

\bibitem{Uribe}
B.~Uribe, \emph{On the classification of pointed fusion categories up to weak Morita equivalence},
Pacific J.\ Math.\ \textbf{290} (2017), 437--466; arXiv:1511.05522.

\bibitem{WG}
Q.-R.~Wang and Z.-C.~Gu, \emph{Towards a complete classification of symmetry-protected topological
phases for interacting fermions in three dimensions and a general group supercohomology theory}, Phys.\
Rev.\ X \textbf{8} (2018), 011055; arXiv:1703.10937.

\end{thebibliography}

\end{document}
